\chapter{Building on this work by yourself}\label{ch:building}

\begin{idea}
Every new analysis in this project follows one recipe: define the unit, build one row per unit with SQL, choose a comparison that removes the obvious confounder, attach an honest uncertainty, and write down what the result does \emph{not} show.
\end{idea}

\section{The recipe in six steps}
\begin{enumerate}
  \item \textbf{Question in one sentence.} ``Is $X$ different between group $G_1$ and group $G_0$?'' If you cannot write it, you are not ready to code.
  \item \textbf{Unit of analysis.} A comment, a pair, a user? The unit decides the SQL and which things are independent (and therefore what to resample).
  \item \textbf{Table of units.} One row per unit, written by a DuckDB query. Print 10 rows and check one by hand against the raw data.
  \item \textbf{Comparison.} Raw difference? Stratified (same comment, same article, same cohort)? Against a null? Name the confounder you are removing.
  \item \textbf{Uncertainty.} Resample the unit that is independent (user, comment). Report an interval, not a $p$-value.
  \item \textbf{Limits.} One paragraph in your notes: what could produce the same result without the effect you think you found?
\end{enumerate}

\section{Starting a new experiment}
\begin{lstlisting}
code/experiments/02-<short-name>/
    common.py        # copy from 01; change nothing but names
    01_<step>.py     # one question per script, with a docstring
notes/experiments/02-<short-name>/
    PLAN.md          # question, units, comparison, planned checks
    NOTES.md         # numbered decisions with [me]/[AI] tags
    RESULTS.md       # every number, with the script that made it
    results/
\end{lstlisting}
Write the \file{PLAN.md} \emph{before} the code, and list the checks you will run (for example, ``also by period, also by topic''). Writing the plan first makes the result honest: a check chosen after seeing the result is not a check.

\section{DuckDB patterns used here}
\begin{lstlisting}
-- 1. Always read Parquet directly; never load a 400M-row table into pandas
SELECT voter_id, count(*) FROM 'derived/votes_clean/*/*.parquet' GROUP BY 1;

-- 2. One 2x2 table per comment: conditional counts with FILTER
SELECT comment_id,
  count(*) FILTER (WHERE replied AND sign = -1) AS a,
  count(*) FILTER (WHERE replied AND sign =  1) AS b,
  count(*) FILTER (WHERE NOT replied AND sign = -1) AS c,
  count(*) FILTER (WHERE NOT replied AND sign =  1) AS d
FROM v GROUP BY 1;

-- 3. Keep only rows that have a match without duplicating columns
SELECT * FROM v SEMI JOIN (SELECT DISTINCT cid FROM repl) p ON v.comment_id = p.cid;
\end{lstlisting}
The script \file{h1\_voter\_strata.py} used in the review (stored in the session scratchpad, not in the repository) is a good small example of pattern 2. If you want it kept, copy it into \file{code/experiments/01-reply-opposition-lifetime/} with a number.

\section{Adding a robustness specification}
Open \file{07\_multiverse.py}. Specifications are built by a nested loop over a small grid (prior $\times$ thresholds $\times$ minimum votes $\times$ include/exclude). To add one dimension, add a list to the grid, add it to the loop and to the output table's columns. Rerun only \file{07}. The figure sorts specifications automatically.

\section{Checking a number you do not trust}
\begin{enumerate}
  \item Find the table in \file{results/} and the line in \file{run.log}.
  \item Recompute it from the Parquet files with a 10-line DuckDB query.
  \item Compute it on 100 rows by hand.
  \item If they differ, the SQL and not the statistics is the usual culprit: check join keys, filters, NULL handling.
\end{enumerate}

\section{Ideas that follow naturally (and how big they are)}
\begin{description}
  \item[Bridging users.] Who replies across the fault line, and are they treated differently? A separate experiment (your decision). It reuses factions (H3) and the reply-back table, plus a definition of ``bridge'' to be pre-registered.
  \item[Reply-choice model.] Fit the popularity--recency--root model of G\'omez et al.\ (2013), with Lan Ceqin's reciprocity extension, per period; use it as a null. Start from her sample design (size-stratified, about 684 threads per period).
  \item[Burst accounts and early prediction.] Add posts per day and the largest number of posts in one thread to the early features, and test how well the first 10 comments predict leaving, following Cheng et al.\ (2015).
  \item[Text.] With embeddings (not downloaded, 508\,GB) one could ask whether opposed replies are semantically distant. Out of reach for now; ask about cluster access.
  \item[Dynamics.] Does a user's stance change after being downvoted? Needs timed votes, which the data lack at fine resolution; month granularity may suffice.
  \item[Causality.] A natural-experiment angle: platform changes in 2022. Needs investigation of the drop first.
\end{description}

\section{Pre-specification, in plain words}
Before running a new analysis, write down in the plan: what you will compute, which outcomes, which subgroups, which thresholds and which checks. Afterwards, any additional analysis is labelled ``exploratory''. This protects you from fooling yourself, and it makes the multiverse chapter of the thesis easy to write. Three lines per analysis are enough.

\section{Writing the methods chapter}
For each analysis, answer: (1) the question; (2) the unit; (3) the exact estimator, with a citation; (4) how uncertainty is computed; (5) the choices that were not forced (pull them from the choice table in Section~\ref{sec:choices}); (6) a limitations sentence. If an answer takes more than a paragraph, split the analysis.

\section{Working habits that save time}
\begin{itemize}
  \item Commit after every working state; a first commit has not been made yet (all files are untracked). A commit message that says what changed and why is documentation.
  \item Never edit derived Parquet files by hand.
  \item Keep AI use labelled: every decision tagged [me], [AI, reviewed] or [AI, not yet reviewed]. Move items to ``reviewed'' only when you could explain them from understanding (\file{notes/method/AI-USE.md}).
  \item When something surprises you, put it in \file{PREGUNTAS.md}; surprises are the raw material for the discussion section.
\end{itemize}
